% Generated draft: requires mathematical and source review.
\documentclass[11pt]{article}
\usepackage[margin=25mm]{geometry}
\usepackage{fontspec}
\setmainfont{Latin Modern Roman}
\usepackage{amsmath,amssymb,mathtools,unicode-math}
\setmathfont{Latin Modern Math}
\usepackage{enumitem,needspace,xurl,hyperref}
\setlength{\parindent}{0pt}
\setlength{\parskip}{5pt}
\setlength{\emergencystretch}{3em}
\newcommand{\N}{\mathbb{N}}
\newcommand{\Z}{\mathbb{Z}}
\newcommand{\Q}{\mathbb{Q}}
\newcommand{\R}{\mathbb{R}}
\newcommand{\C}{\mathbb{C}}
\newcommand{\F}{\mathbb{F}}
\newcommand{\A}{\mathbb{A}}
\newcommand{\PP}{\mathbb{P}}
\newcommand{\E}{\mathbb{E}}
\newcommand{\eps}{\varepsilon}
\newcommand{\dd}{\,\mathrm{d}}
\begin{document}
\section*{UnsolvedMath ID: 30004032}
\textbf{Source problem code:} catalogue entry
\subsection{Definitions and mathematical statement}
A plane graph is a graph with a fixed embedding in the plane. A near-triangulation has triangular bounded faces. The graphs considered in the supplied statement have outer face bounded by a cycle of length five. The record does not specify further graph conventions, such as simplicity or connectedness.

The supplied conjecture is that, for every such graph \(G\),
\[
 n_G^\star=\sum_{i=1}^{11}\lambda_i n_{G_i}
 \qquad\text{for some }(\lambda_1,\ldots,\lambda_{11})\in\R_{\geq 0}^{11}.
\]
The coefficients may depend on \(G\), whereas the eleven generators are intended to be fixed across the assertion.

\textbf{Unresolved definitions.} The supplied record does not define \(n_G\), the operation denoted by the superscript \(\star\), the eleven graphs \(G_1,\ldots,G_{11}\), or the common vector space in which the equality is to hold. Thus the displayed formula records the asserted cone membership, but does not give a self-contained mathematical conjecture until these data are supplied. No particular coloring interpretation or choice of generators is assumed here.

The original statement reports computational verification for graphs with at most \(26\) faces. This is a supplied computational claim, not a proof of the universally quantified assertion. The record does not specify whether the outer face is included in that count or how the computation was performed.
\subsection{Short English statement}
The conjecture asks whether every plane near-triangulation with a five-cycle as its outer boundary has its quantity \(n_G^\star\) in the nonnegative cone generated by eleven fixed quantities \(n_{G_i}\). The quantities, the superscript operation, and the eleven generators are missing from the supplied record, so the formulation remains schematic and requires human mathematical review.
\subsection{Sources}
\begin{itemize}
\item Supplied local UnsolvedMath registry record, permanent ID \(30004032\), source code \texttt{catalogue entry}; registry locations specified by the project are \texttt{data/unsolvedmath/problems.json} and the supplementary \texttt{catalog/problems.json}. Title: Five-Cycle Near-Triangulation Decomposition Conjecture. The record gives proposed year \(2019\) and does not name a proposer.
\item Supplied report citation: \emph{Graph Theory (2020)}, Oberwolfach Reports, DOI 10.4171/owr/2019/1, pp.\ 5--63. \url{https://doi.org/10.4171/owr/2019/1}. The background describes Oberwolfach Reports as documenting mathematical workshops and their open-problem discussions; no precise location of this conjecture within the report is supplied.
\item Dvorak and Lidicky, \emph{Coloring count cones}: \url{https://arxiv.org/abs/1907.04066}.
\item Supplied Journal of Graph Theory version: \url{https://onlinelibrary.wiley.com/doi/10.1002/jgt.22767}.
\item Literature assessment dated exactly \texttt{2026-08-22}: the supplied curation labels the problem open and reports finite computational verification and later obstruction results, with no general proof or counterexample resolving the conjecture. This is supplied dated curation derived from the OWR record, not an independently verified literature assessment; its status and cited claims require independent verification.
\end{itemize}
\end{document}
